% Part: first-order-logic
% Chapter: tableaux

\documentclass[../../../include/open-logic-chapter]{subfiles}

\begin{document}

\iftag{FOL}
      {\olchapter{fol}{tab}{ট্যাবলো}}
      {\olchapter{pl}{tab}{ট্যাবলো}}

\begin{editorial}
এই অধ্যায়ে একটি চিহ্নিত বিশ্লেষণাত্মক ট্যাবলো পদ্ধতি উপস্থাপন করা হয়েছে।

প্রমাণ-পদ্ধতি হিসেবে ট্যাবলো-সংক্রান্ত উপাদান অন্তর্ভুক্ত বা বাদ দিতে
“prfTab” ট্যাগটি ব্যবহার করো।
% BN-SRC-043 TARGET-CORRECTION-BEGIN
\par\smallskip\noindent\textnormal{\small\textbf{উৎস-সংশোধন (BN-SRC-043):} হিমায়িত ইংরেজি গদ্যে “prfTab” ট্যাগটিকে ভুল করে স্বাভাবিক নিষ্পাদন-সংক্রান্ত উপাদানের নিয়ন্ত্রক বলা হয়েছে। ট্যাগের নাম, অধ্যায়ের বিষয় এবং ব্যবহারের সঙ্গে মিলিয়ে বাংলা গদ্যে ট্যাবলো বলা হয়েছে।} % BN-SRC-043 TARGET-CORRECTION-END
\end{editorial}

\olimport{rules-and-proofs}

\olimport{propositional-rules}

\iftag{FOL}{%
  \olimport{quantifier-rules}
}{}

\olimport{derivations}

\olimport{proving-things}

\iftag{FOL}{%
  \olimport{proving-things-quant}
}{}

\olimport{proof-theoretic-notions}

\olimport{provability-consistency}

\olimport{provability-propositional}

\iftag{FOL}{%
  \olimport{provability-quantifiers}
}{}

\olimport{soundness}

\iftag{FOL}{%
  \olimport{identity}
  \olimport{soundness-identity}
}{}

\OLEndChapterHook

\end{document}
